% Authors: Dr. Denys Dutykh (Mathematics Department, Khalifa University of Science and Technology,
% Abu Dhabi, UAE) and Prof. Laurent Vuillon (Univ. Savoie Mont Blanc, CNRS, LAMA, Chambery, France)
\section{Occurrences, representations, and Euclidean cost}\label{sec:occurrences}

The arithmetic comparison in this paper concerns every primitive
representation of a fixed integer, including representations that do not
come from a Markov occurrence. We define the cost, put the two sets of
objects on the same cost scale, and explain why minimality of the cost is
enough to separate two occurrences.

\subsection{The Euclidean cost and its two readings}\label{sec:cost}

The cost of a pair is read in two ways in this paper: representations are
compared through their continued fractions, while the cost of an occurrence
is computed from a product of matrices. We define the cost by subtractions
and prove that both readings give it. For coprime positive integers $r,s$,
let $\ell(r,s)$ be the number of single subtractions needed to reach
$(1,1)$: subtract the smaller coordinate from the larger until equality. In particular,
$\ell(1,1)=0$, and exchanging the coordinates does not change $\ell$. For a
column $v$ with coprime positive entries we write $\ell(v)$ for the cost of
its two entries. Define
\begin{equation}\label{eq:continued-fraction-matrices}
 P_j=\begin{pmatrix}j&1\\1&0\end{pmatrix}\quad(j\geqslant1).
\end{equation}
For every primitive positive column $v$, exactly $j$ initial
subtractions carry $P_jv=(jv_1+v_2,v_1)^{\mathsf T}$ to
$(v_2,v_1)^{\mathsf T}$. Each of those subtractions is permitted,
including the final step when $v_1=v_2=1$. Consequently
\begin{equation}\label{eq:digit-adds-cost}
 \ell(P_jv)=j+\ell(v).
\end{equation}

\begin{lemma}[The cost by continued fractions and by matrices]\label{lem:cost-descriptions}
Let $q_0,q_1,\ldots,q_j$ be positive integers, with $q_j\geqslant2$ unless
$j=0$ and $q_0=1$. Then the column
\begin{equation}\label{eq:continued-fraction-product}
 \binom sr=P_{q_0}P_{q_1}\cdots P_{q_j}\binom10
\end{equation}
has coprime positive entries, $s/r=[q_0;q_1,\ldots,q_j]$, and
\begin{equation}\label{eq:quotient-cost}
 \ell(r,s)=q_0+q_1+\cdots+q_j-1 .
\end{equation}
Every pair of coprime integers $0<r\leqslant s$ arises in this way, from
the canonical continued fraction of $s/r$. If $(r,s)\neq(1,1)$, then
moreover
\begin{equation}\label{eq:index-product}
 \binom sr=P_{q_0}\cdots P_{q_{j-1}}P_{q_j-1}\binom11 ,
\end{equation}
and $\ell(r,s)$ is the sum of the indices of the matrices in this product.
\end{lemma}

\begin{proof}
\emph{The product.} We argue by induction on $j$. If $j=0$, the column is
$P_{q_0}(1,0)^{\mathsf T}=(q_0,1)^{\mathsf T}$, and $s/r=q_0=[q_0]$. If $j\geqslant1$,
the induction hypothesis applied to $q_1,\ldots,q_j$ gives a column
$(s',r')^{\mathsf T}=P_{q_1}\cdots P_{q_j}(1,0)^{\mathsf T}$ with positive entries and
$s'/r'=[q_1;q_2,\ldots,q_j]$, so that
$$
 \binom sr=P_{q_0}\binom{s'}{r'}=\binom{q_0s'+r'}{s'},\qquad
 \frac sr=q_0+\frac{r'}{s'}=q_0+\frac1{[q_1;q_2,\ldots,q_j]}=[q_0;q_1,\ldots,q_j].
$$
In both cases the entries are positive, and they are coprime because each
matrix $P_q$ has determinant $-1$, so that the column is the image of the
primitive column $(1,0)^{\mathsf T}$ under a matrix invertible over $\Z$.

\emph{Every pair arises.} Let $0<r\leqslant s$ be coprime. The canonical
continued fraction $[q_0;q_1,\ldots,q_j]$ of $s/r$ has positive quotients, since
$s/r\geqslant1$, and its last quotient is at least $2$ unless $s/r=1=[1]$. The
column built from it has coprime positive entries in the ratio $s/r$, and two
primitive positive columns with the same ratio are equal; hence it is
$(s,r)^{\mathsf T}$.

\emph{The cost.} For $0\leqslant i\leqslant j$ put
$v_i=P_{q_i}\cdots P_{q_j}(1,0)^{\mathsf T}$, so that $v_0=(s,r)^{\mathsf T}$,
$v_j=(q_j,1)^{\mathsf T}$ and $v_i=P_{q_i}v_{i+1}$ for $i<j$. By the first step,
every $v_i$ is a primitive positive column, so \eqref{eq:digit-adds-cost},
applied $j$ times, gives
$$
 \ell(r,s)=q_0+q_1+\cdots+q_{j-1}+\ell(q_j,1).
$$
From the pair $(q_j,1)$ every subtraction lowers the first coordinate by one,
and $(1,1)$ is reached after $q_j-1$ of them, so that $\ell(q_j,1)=q_j-1$. This
proves \eqref{eq:quotient-cost}.

\emph{The second product.} If $(r,s)\neq(1,1)$, then $q_j\geqslant2$, so that the
matrix $P_{q_j-1}$ is defined, and
$P_{q_j-1}(1,1)^{\mathsf T}=(q_j,1)^{\mathsf T}=P_{q_j}(1,0)^{\mathsf T}$. Substituting
this in \eqref{eq:continued-fraction-product} gives \eqref{eq:index-product}.
The indices of its factors add up to $q_0+\cdots+q_{j-1}+(q_j-1)$, which is
$\ell(r,s)$ by \eqref{eq:quotient-cost}.
\end{proof}

The introduction described the cost through the continued fraction, as in
\eqref{eq:quotient-cost}. The proof of \cref{prop:occurrence-cost} below
computes it through products of the matrices $P_1$ and $P_2$, by
\eqref{eq:digit-adds-cost}, as the sum of the indices in
\eqref{eq:index-product}. \Cref{lem:cost-descriptions} shows that the two
readings agree, and \cref{lem:euclidean-bijection} adds a third, the
length of a word.

\subsection{The occurrence-induced pair}

For $0<P<Q$ with $\gcd(P,Q)=1$, let $\gamma_{P/Q}$ be the binary
word of length $Q$ whose $i$th letter is $b$ precisely when
$\lfloor iP/Q\rfloor-\lfloor(i-1)P/Q\rfloor=1$.
It is the proper lower Christoffel word with
\begin{equation}\label{eq:farey-counts}
 |\gamma_{P/Q}|_a=Q-P,\qquad |\gamma_{P/Q}|_b=P.
\end{equation}
Thus $Q$ is the total length of the Christoffel word. This convention
distinguishes the denominator from the number of letters $a$.
The word starts with $a$ and ends with $b$, so write
$\gamma_{P/Q}=azb$. Its interior $z$ is a palindrome.
Indeed, for $1\leqslant j\leqslant Q-1$, coprimality gives
$\lfloor(Q-j)P/Q\rfloor=P-1-\lfloor jP/Q\rfloor$.
The increments defining the letters at positions $i$ and $Q+1-i$
therefore agree for $2\leqslant i\leqslant Q-1$.
This is the central-word property of Christoffel words
\cite{BertheDeLucaReutenauer2008}. We write $\widetilde w$ for reversal.

For later exact evaluation, the same occurrence has the Cohn description
\begin{equation}\label{eq:cohn-source}
 \mu(a)=\begin{pmatrix}2&1\\1&1\end{pmatrix},\qquad
 \mu(b)=\begin{pmatrix}5&2\\2&1\end{pmatrix},\qquad
 m_{P/Q}=\frac13\tr\mu(\gamma_{P/Q}).
\end{equation}
The standard Christoffel--Markov correspondence identifies these
occurrences with normalized Markov triples having largest entry greater
than $2$; see, for example,
\cite[Theorem~1 and its proof]{ReutenauerVuillon2017}.
No injectivity of $P/Q\mapsto m_{P/Q}$ is assumed.

The representation attached to an occurrence, and its cost, are known.
\c{C}anak\c{c}{\i} and Schiffler show that $m_{P/Q}$ is the numerator of a
palindromic continued fraction whose entries are $1$ and $2$, with
$2(Q-P-1)$ entries equal to $1$ and $2P$ entries equal to $2$, and that
$m_{P/Q}=r^2+s^2$ where $s/r$ is the continued fraction of its right half
\cite[Section~4.1, Theorem~4.1 and the proof of Corollary~4.4]{CanakciSchiffler2020}.
That half has partial quotients summing to $(Q-P-1)+2P=P+Q-1$ and ending
with $2$, so \cref{lem:cost-descriptions} gives the cost $P+Q-2$. The same statement, in the language of Euclidean paths, is the
source dictionary of our companion article \cite{DutykhVuillon2026Endpoint}.
Alessandro De Luca and the second author have another proof of the
representation, within the combinatorics of words (private communication), in
which the Christoffel word of length $m_{P/Q}$ given by the construction of
Reutenauer and Vuillon \cite{ReutenauerVuillon2017} is the image of a word with
$s$ letters $a$ and $r$ letters $b$ under a substitution that sends $a$ to a
word of length $s$ and $b$ to a word of length $r$, so that its length is
$s^2+r^2$.
We include a proof in the matrix language of the present paper, because it
defines the pair that the computations of \cref{sec:gaussian} use; the
paragraph after the proof shows that it is the pair of \c{C}anak\c{c}{\i} and
Schiffler.

\begin{proposition}[Occurrence pair and its cost]\label{prop:occurrence-cost}
Every occurrence $P/Q$ with $0<P<Q$ induces a primitive positive
representation
$$
 m_{P/Q}=r_{P/Q}^2+s_{P/Q}^2,\qquad 0<r_{P/Q}<s_{P/Q}.
$$
Its subtractive Euclidean cost is $P+Q-2$.
\end{proposition}

\begin{proof}
\emph{The label as a quadratic form.} The matrices $\mu(a)$ and $\mu(b)$
are symmetric, so reversing a word transposes its Cohn matrix, and the
palindrome $z$ has a symmetric matrix
$\mu(z)=\left(\begin{smallmatrix}\alpha&\beta\\\beta&\delta\end{smallmatrix}\right)$.
Multiplying out,
$$
 \mu(a)\,\mu(z)\,\mu(b)=
 \begin{pmatrix}
 10\alpha+9\beta+2\delta & 4\alpha+4\beta+\delta\\
 5\alpha+7\beta+2\delta & 2\alpha+3\beta+\delta
 \end{pmatrix},
$$
whose trace is $3(4\alpha+4\beta+\delta)$. With $e=(2,1)^{\mathsf T}$ one has
$e^{\mathsf T}\mu(z)e=4\alpha+4\beta+\delta$, and therefore
\begin{equation}\label{eq:square-count-identity}
 m_{P/Q}=\tfrac13\tr\mu(\gamma_{P/Q})=e^{\mathsf T}\mu(z)\,e .
\end{equation}

\emph{A square root of the form.} Cut the palindrome at its middle:
$z=\widetilde w\,w$ when $\abs z$ is even, and $z=\widetilde w\,a\,w$ or
$z=\widetilde w\,b\,w$ when $\abs z$ is odd, according to its middle letter.
Since $\mu(\widetilde w)=\mu(w)^{\mathsf T}$, $\mu(a)=P_1^2$ and
$\mu(b)=P_2^2$, with $P_1$ and $P_2$ symmetric, we obtain
$\mu(z)=H^{\mathsf T}H$, where $H$ is $\mu(w)$, $P_1\mu(w)$ or $P_2\mu(w)$ in
the three cases; the empty palindrome gives $H=I$. By
\eqref{eq:square-count-identity}, $m_{P/Q}=(He)^{\mathsf T}(He)$.

\emph{The pair.} The integer matrix $H$ has determinant $\pm1$, so the
column $He$ is primitive. Its entries are positive, because $H$ is a product
of nonnegative matrices with no zero row and $e$ is positive. They are
distinct: two equal coprime entries would both be $1$, giving
$m_{P/Q}=2$, whereas $m_{P/Q}\geqslant5$, because $\mu(z)\geqslant I$ entrywise
(a product of nonnegative matrices each at least $I$) and
\eqref{eq:square-count-identity} then gives $m_{P/Q}\geqslant e^{\mathsf T}e=5$. Sorting the two entries defines
$r_{P/Q}<s_{P/Q}$. The construction uses the occurrence alone and assumes
nothing about other occurrences with the same label.

\emph{Its cost.} The column $e$ reaches $(1,1)$ after one subtraction, so
$\ell(e)=1$. Apply the factors of $H$ to $e$ from the right. By
\eqref{eq:digit-adds-cost}, a factor $\mu(a)=P_1P_1$ coming from a letter
$a$ of $w$ adds $2$ to the cost, a factor $\mu(b)=P_2P_2$ adds $4$, and the
middle factor $P_1$ or $P_2$, when present, adds $1$ or $2$. Each letter of
$w$ occurs twice in $z$, so these increments add up to $\abs z_a+2\abs z_b$,
and \eqref{eq:farey-counts} gives
$$
 \ell(He)=1+|z|_a+2|z|_b
 =1+(Q-P-1)+2(P-1)=P+Q-2.
$$
No restriction on the parity of $|z|$ or of $m_{P/Q}$ is needed.
\end{proof}

The pair of \c{C}anak\c{c}{\i} and Schiffler is this pair. Write
$\gamma_{P/Q}=a^{u_1}b\,a^{u_2}b\cdots a^{u_P}b$. The $k$th letter $b$ stands at
position $\lceil kQ/P\rceil$, and $kQ/P$ is not an integer for $1\leqslant k<P$;
comparison with their crossing numbers $v_k$, which satisfy
$v_1+\cdots+v_k=\lfloor kQ/P\rfloor$ for $k<P$ and $v_1+\cdots+v_P=Q-1$, gives
$u_1=v_1$ and $u_k=v_k-1$ for $2\leqslant k\leqslant P$. Their palindromic
continued fraction is therefore $[2,\sigma(z),2]$, where $\sigma$ replaces each
letter $a$ by the two entries $1,1$ and each letter $b$ by $2,2$. Its right
half is $[\sigma(w),2]$, or $[j,\sigma(w),2]$ with $j=1$ or $2$ for a middle
letter $a$ or $b$, and since $\mu(a)=P_1^2$, $\mu(b)=P_2^2$ and $e=P_2(1,0)^{\mathsf T}$,
this is the continued fraction of the column $He$.

\subsection{Every primitive pair has a unique Euclidean directive}

Introduce the elementary matrices
\begin{equation}\label{eq:directive-incidence}
 L_a=\begin{pmatrix}1&1\\0&1\end{pmatrix},\qquad
 L_b=\begin{pmatrix}1&0\\1&1\end{pmatrix},\qquad
 L_{uv}=L_uL_v,\qquad {\bf1}=\binom11.
\end{equation}

\begin{lemma}[Euclidean bijection]\label{lem:euclidean-bijection}
The map $u\mapsto L_u{\bf1}$ is a bijection from binary words to
ordered coprime positive pairs, and
\begin{equation}\label{eq:euclidean-length}
 \ell(L_u{\bf1})=|u|.
\end{equation}
\end{lemma}

\begin{proof}
If $x>y>0$, the first letter of any directive producing $(x,y)$ must
be $a$, because $L_a$ maps positive pairs to pairs with first coordinate
larger, whereas $L_b$ makes the second coordinate larger. The unique
predecessor is $(x-y,y)$. For $y>x$, the first letter is $b$ and the
predecessor is $(x,y-x)$. Each subtraction preserves positivity and
coprimality and decreases $x+y$. The only equal coprime positive pair
is $(1,1)$, so this process terminates there. Reading the uniquely
determined letters gives both existence and uniqueness, and one letter
is removed at each subtraction. This proves
\eqref{eq:euclidean-length}.
\end{proof}

The word $u$ is the \emph{directive} of the pair $L_u{\bf1}$, and it can be
read on the continued fraction. With
$J=\left(\begin{smallmatrix}0&1\\1&0\end{smallmatrix}\right)$, which exchanges
the two coordinates, one has $P_q=L_a^qJ$, $JL_a^qJ=L_b^q$ and $J{\bf1}={\bf1}$.
Grouping the factors of \eqref{eq:index-product} in consecutive pairs
therefore gives, for coprime $0<r<s$ with $s/r=[q_0;q_1,\ldots,q_j]$,
$$
 \binom sr=L_a^{q_0}L_b^{q_1}L_a^{q_2}L_b^{q_3}\cdots\,{\bf1},
$$
where the exponents are $q_0,q_1,\ldots,q_{j-1},q_j-1$ and the letters
alternate. The directive of $(s,r)$ is thus the word
$a^{q_0}b^{q_1}a^{q_2}\cdots$ whose runs have these lengths, and its length
$q_0+\cdots+q_j-1$ agrees with \eqref{eq:quotient-cost} and
\eqref{eq:euclidean-length}. For instance $21/13=[1;1,1,1,1,2]$, and the
directive of $(21,13)$ is $ababab$.

For unordered representations we always choose $0<r<s$. This loses
neither norms nor costs. The exceptional positive pair $(1,1)$ represents
$2$ and can be treated separately; $1=1^2+0^2$ is outside the positive
pair convention.

\begin{definition}\label{def:weak-minimality}
An occurrence $P/Q$ is \emph{weakly Euclidean-minimal} if
\begin{equation}\label{eq:weak-minimality}
 P+Q-2\leqslant\ell(r,s)
 \quad\text{whenever}\quad
 m_{P/Q}=r^2+s^2,\quad 0<r<s,\quad\gcd(r,s)=1.
\end{equation}
\end{definition}

The competitors in \eqref{eq:weak-minimality} range over the entire
primitive two-square fiber. Their Euclidean directives need not be
directives induced by a central Cohn word as above.
Ties with such other representations are allowed.

\subsection{Why weak minimality is enough within a proved family}

Among the ordering theorems recalled in the introduction, we use the
fixed-sum theorem of Lagisquet, Pelantov\'a, Tavenas and Vuillon
\cite[Theorem~2, item~3]{LagisquetEtAl2021}, which compares two lattice points
differing by a unit step along a line of fixed sum.
Applying it $i$ times gives, for reduced $P/Q$ and
$(P-i)/(Q+i)$ with $0<i<P<Q$,
\begin{equation}\label{eq:fixed-sum-order}
 m_{P/Q}<m_{(P-i)/(Q+i)}.
\end{equation}
The intermediate points of that iteration need not be reduced; this is
harmless because the authors define the label at every lattice point,
not only at Farey fractions.
The inequality compares distinct points on one line of fixed sum $P+Q$.
It does not compare arbitrary two-square representations.

The next lemma is an immediate consequence of \eqref{eq:fixed-sum-order};
it is also the case of equal sums in the proof of the proposition on the
lengths of directing words of Markov words in
\cite{DutykhVuillon2026Reflections}, where the cost $P+Q-2$ appears as the
length of a directing word.

\begin{lemma}[Fixed-cost injectivity]\label{lem:fixed-cost-injectivity}
Two occurrences having the same label and the same occurrence cost
have the same reduced Farey coordinate.
\end{lemma}

\begin{proof}
Their costs are $P+Q-2$ and $P'+Q'-2$. If they are equal and
$P>P'$, set $i=P-P'$. Then $Q'=Q+i$ and $0<i<P$, so
\eqref{eq:fixed-sum-order} contradicts equality of labels. The other
strict ordering of the numerators is excluded in the same way.
Consequently $P=P'$ and $Q=Q'$. A reduced Farey coordinate has a unique
lower Christoffel word, hence a unique normalized occurrence.
\end{proof}

\begin{corollary}\label{cor:restricted-injectivity}
The label map is injective on any family of occurrences all of which
are weakly Euclidean-minimal.
\end{corollary}

\begin{proof}
If two members have equal label, each occurrence-induced pair is a
competitor in the other's two-square fiber. Weak minimality in both
directions makes their costs equal. Apply
Lemma~\ref{lem:fixed-cost-injectivity}.
\end{proof}

Thus weak minimality on every occurrence would imply the
Frobenius--Markov uniqueness conjecture. For a restricted family, both occurrences
must lie in that family. If only one is known to be weakly minimal, a
distinct equal-label occurrence is merely forced to have strictly
larger cost.
